%% kapa: exam=1 round=22 year=2011 period=1 form=A subject=economics subject_ko=경제원론 seq=14 number=54
%% source: Q-Net 감정평가사 기출문제 (cst00309) artlSeq=5036790 file=extracted/1008/5036790/2011년 제22회 감정평가사 1차시험 기출문제/1교시.hwp
%% conversion: kordoc 4.24.51 -> tools/gen1.py (auto); review=proofread-sonnet; flags=none
\begin{question}{54}
생산함수가 Q=5L\textsuperscript{0.4}K\textsuperscript{0.6} 일 때, 다음 설명 중 옳은 것은? (단, Q, L, K는 각각 생산량, 노동 투입량, 자본 투입량, Q>0, L>0, K>0)
\begin{choices}
  \item L=K일 경우 노동의 한계생산은 일정하다.
  \item 노동과 자본간의 대체탄력성은 L, K 값의 크기에 따라 변한다.
  \item 등량곡선은 우하향하는 직선 모양을 갖는다.
  \item 규모에 대한 수익이 체감한다.
  \item 한계기술대체율은 L, K 값의 크기와 관계없이 항상 일정하다.
\end{choices}
\end{question}
\answer{1}
